\chapter{Materiais e Metodologia}

A presente pesquisa caracteriza-se como uma pesquisa aplicada, com abordagem mista (qualitativa e quantitativa) e objetivos descritivos. O método fundamentou-se na análise técnica da integração entre os princípios de DevSecOps e a Infraestrutura como Código (IaC), culminando na construção de uma Prova de Conceito (PoC) para validação do modelo proposto.

\section{Classificação da Pesquisa}
Quanto à sua natureza, esta pesquisa é classificada como aplicada, pois objetiva gerar conhecimentos para aplicação prática e solução de problemas específicos de segurança em infraestrutura. Do ponto de vista dos objetivos, é descritiva e exploratória, ao estabelecer relações entre ferramentas de automação e conformidade. A abordagem quantitativa foi empregada na coleta de métricas de desempenho e eficácia da esteira automatizada.

\section{Etapas do Procedimento Metodológico}
A execução do trabalho foi estruturada em quatro etapas lógicas, detalhadas a seguir:


\begin{enumerate}
    \item \textbf{Revisão Bibliográfica:}
    
    Realização de levantamento sistemático na literatura sobre o estado da arte de DevSecOps, IaC e ferramentas de análise estática (SAST). Foram consultados repositórios acadêmicos, documentações oficiais de provedores de nuvem (AWS) e frameworks de mercado (CIS Benchmarks).
    
    \item \textbf{Definição do Cenário e Seleção de Ferramentas:}
    
    Identificação dos recursos críticos de infraestrutura (EC2, S3, Redes) e seleção do stack tecnológico líder de mercado capaz de suportar a automação de segurança.
    
    \item \textbf{Desenvolvimento da Prova de Conceito (PoC):}
    
    Construção da infraestrutura utilizando Terraform e integração da esteira de CI/CD. Nesta fase, foram configuradas as políticas de segurança e os gatilhos de interrupção (fail-fast).
    
    \item \textbf{Validação do Modelo:}
    
    O modelo proposto será submetido à avaliação de especialistas em segurança, DevOps e IaC, por meio de workshops ou entrevistas, para obter feedback e identificar possíveis melhorias.
    
    \item \textbf{Coleta e Análise de Dados:}
    
    Submissão do código a testes de vulnerabilidade propositais para validar a capacidade de detecção das ferramentas e mensurar o tempo de resposta em comparação a processos manuais.
    
\end{enumerate}


\section{Materiais e Ecossistema Tecnológico}
Para a viabilização dos testes práticos e validação do modelo, utilizou-se um conjunto de recursos de software e hardware que compõem o ecossistema de desenvolvimento moderno:

\begin{itemize}
    \item\textbf{Ambiente de Desenvolvimento:} Estação de trabalho com sistema operacional Linux (Ubuntu 22.04 via WSL2) e ambiente de desenvolvimento integrado (IDE) VS Code.
    
    \item\textbf{Provedor de Nuvem:} Amazon Web Services (AWS) como ambiente de destino para o provisionamento dos recursos de computação e armazenamento.
    
    \item\textbf{Orquestração e IaC:} Terraform para a definição declarativa da infraestrutura e Ansible para gerenciamento de configurações.

    \item \textbf{Ferramental de Segurança:}
    A pesquisa validou indicadores fundamentais para medir o sucesso da integração:
    \begin{itemize}
        \item \textbf{Checkov:} Motor de análise estática para detecção de falhas de configuração em arquivos Terraform.
        \item \textbf{GitGuardian:} Ferramenta dedicada à detecção e prevenção de vazamento de segredos (secrets scanning).
        \item \textbf{Open Policy Agent (OPA):} Motor de políticas utilizado para aplicar regras de governança e FinOps via linguagem Rego.
    \end{itemize}

     \item\textbf{Controle de Versão e CI/CD:} Repositórios Git hospedados em instâncias de GitLab/GitHub, responsáveis pelo acionamento automático dos pipelines de validação.
\end{itemize}

\section{Procedimentos de Validação}
A validação do modelo ocorreu por meio da execução de test-cases onde configurações inseguras (ex: portas SSH abertas para a internet e buckets S3 públicos) foram codificadas intencionalmente. O sucesso da metodologia foi medido pela capacidade da esteira em identificar essas falhas em tempo real (fase de build) e impedir o provisionamento de recursos não conformes, garantindo a integridade do ambiente.